\chapter{La música árabe}

    La música árabe tiene unas características propias que es necesario conocer para así entender la manera correcta de escucharla e interpretarla. En este capítulo, veremos las ideas fundamentales que la definen.

    \section{¿Qué es la música árabe?}

    Bajo el término <<música árabe>> se engloban las expresiones musicales que aparecen, con una serie de características comunes, en los países del mundo árabe. Estas características son el resultado de procesos tanto culturales como políticos, y de influencias a lo largo de su historia, entre las que cabe destacar:

    \begin{itemize}
    \item Influencia de las culturas limítrofes durante los primeros siglos del Islam. Entre ellas, las de Bizancio, Siria, Mesopotamia y Persia, cada una con sus propias tradiciones musicales.
    \item Influencia del mundo clásico. En especial de Grecia, que dio como resultado la escritura de tratados sobre la música y una cierta formalización del hecho músical árabe.
    \item Influencia occidental durante el Medievo. Los desarrollos artísticos que tuvieron lugar en Al-Ándalus llegaron después al magreb tras la expulsión de los moriscos, para continuar allí su evolución. La figura de Ziryab, ya mencionada al tratar la historia del ud, es de especial relevancia también aquí.
    \item Influencia otomana. La hegemonía de los turcos otomanos entre los siglos \textsc{XV} y \textsc{XX} tuvo gran influencia en la música árabe, que asimiló elementos de la música turca, la cual a su vez había recibido antes la influencia de las de Asia Central, Anatolia o Iraq, entre otras.
    \item Influencia occidental moderna. En especial desde la conquista de Egipto por parte de Napoleón. La música árabe recibe elementos de la música occidental, que van desde formas musicales a la propia notación empleada hoy día.
    \end{itemize}

    Aunque la música árabe es muy variada y existen muchas formas regionales, lo que hoy día conocemos como música árabe es, fundamentalmente, la música que corresponde al periodo denominado <<Edad Dorada de la Música Árabe>>, que es en realidad una época de esplendor de la música egipcia. A lo largo de un periodo de 40 años centrados alrededor de la mitad del siglo \textsc{xx}, El Cairo se convirtió en la capital músical del mundo árabe, y durante este tiempo se produjo una evolución musical sin precedentes. Compositores como Sayed Darwish, Mohammad Al Qasabgi, Riad Al Sunbati o Mohammad Abdel Wahab fueron los principales impulsores de un cambio profundo en la música árabe. 

    Siendo esta la época en la que el mercado discográfico comenzaba a instaurarse, y coincidiendo a su vez con un periodo también de esplendor del cine egipcio, las nuevas formas musicales se difundieron ampliamente por todo el mundo árabe. Gracias a cantantes como Umm Kalzum o Layla Murad, de inmensa popularidad, las composiciones y estilos de estos compositores se establecieron como parte fundamental del repertorio de la música árabe contemporánea.

    \section{Características}

    Los principales aspectos que definen la música árabe son los siguientes. Algunos de ellos son comunes a los de otras tradiciones musicales, mientras que otros son específicos.

    \begin{itemize}
    \item \textbf{Carácter eminentemente vocal}. La voz es el instrumento principal de la música árabe, y el cantante o la cantante tiene una entidad muy superior a la del instrumentista. La preeminencia de la voz frente a los instrumentos puede verse en la propia música como tal, supeditada a las particularidades de esta. Incluso en las piezas instrumentales, y especialmente en las improvisadas, la forma canónica de tocar un instrumento busca imitar a la voz. El rango tonal más frecuente, el desarrollo melódico o el uso de silencios y pausas, entre otras propiedades, están encaminadas a conseguir con el instrumento una sonoridad similar a la de la voz, y a la manera en que esta emplea la respiración o los melismas que son característicos de la improvisación vocal.

    \item \textbf{Importancia de la melodía y ausencia de armonía}. A diferencia de la música occidental, la armonía no existe en la música árabe. La melodía y el ritmo son los pilares de una pieza, y es sobre ellos que se hace énfasis. El desarrollo melódico, así como la riqueza expresiva que el intérprete añade mediante la ornamentación (ver capítulo \ref{ornamentacion}), son los elementos que aportan interés y valor a la melodía, sin necesidad de un contexto armónico ni de las tensiones que este implica.

    El acompañamiento se limita a la interpretación de la misma línea melódica con ciertas variaciones, o bien a la presencia de una nota pedal (frecuentemente la tónica del modo melódico empleado).

    Aunque algunas regiones tienen formas musicales polifónicas, se puede considerar que la música árabe es de carácter monofónico.

    \item \textbf{El sistema tonal árabe}. El sistema tonal árabe comporta no solo una subdivisión de la octava en intervalos distintos a los empleados en la música occidental, sino también una manera propia de utilizar los distintos tonos y construir melodías.

    \item \textbf{La instrumentación}. La música árabe se interpreta con unos instrumentos particulares y unas técnicas específicas. También la composición de los conjuntos musicales tiene formas propias, y todo ello caracteriza la manera en la que se interpreta el repertorio.

    La ausencia de armonía, por ejemplo, se refleja en los instrumentos típicos árabes, que son de tipo melódico, y con los que no es posible (o resulta muy complejo) el tocar acordes.
    \end{itemize}

\label{sistema_tonal}
\section{Sistema tonal}

    Mientras que la música occidental se basa en las 12 notas de la escala cromática (12 subdivisiones de una octava), en la música árabe la octava se divide en 24 notas. Es decir, aparece una nueva nota entre cada dos notas de la escala cromática, y la separación entre dos notas consecutivas no es de un semitono, sino de un cuarto de tono.

    A efectos de notación, se utilizan los conceptos de \emph{semibemol} y \emph{semisostenido}, alteraciones para las que se adoptan los símbolos \scalerel*{\includegraphics{semibemol.png}}{B} y \scalerel*{\includegraphics{semisostenido.png}}{B}. En la tablatura que se emplea en este libro, las posiciones correspondientes a estas notas intermedias situadas entre las de la escala cromática se indican añadiendo $\frac12$ a la cifra de la posición que representa la nota natural. Por ejemplo:

    \begin{center}
    \begin{minipage}{.2\textwidth}
        \centering
        \begin{lilypond}
            \include "arabic.ly"

            notes = {
                \cadenzaOn
                mib misb mi
                \bar "||"
            }

            \layout {
              indent = #0
              \context {
                \Score
                supportNonIntegerFret = ##t
              }
              \context {
                  \TabStaff
                stringTunings = \stringTuning <do, fa, la, re sol do'>
                }
              \context {
                  \Staff
                  \remove "Time_signature_engraver"
                }
            }

            <<
              \new Staff << \clef "G_8" \notes >>
              \new TabStaff << \notes>>
            >>
        \end{lilypond}
    \end{minipage}
    \hfill
    \begin{minipage}{.79\textwidth}
        \begin{flushright}
        \resizebox{.9\textwidth}{!}{
        \begin{tikzpicture}[
            ynode/.style={scale=.35,inner sep=1pt}]
          \fretboard

        \draw[blue!50, fill=blue!50] (4-2) circle [radius=0.09] node[scale=.30] {};
        \node[ynode] at (4-2) {\small Mi$\flat$};

        \draw[blue!50, fill=blue!50] (4-3) circle [radius=0.09] node[scale=.30] {};
        \node[ynode] at (4-3) {\small Mi\hflat};

        \draw[blue!50, fill=blue!50] (4-4) circle [radius=0.09] node[scale=.30] {};
        \node[ynode] at (4-4) {\small Mi};
        \end{tikzpicture}
        }
        \end{flushright}
    \end{minipage}
    \end{center}

    En la práctica, de los símbolos de alteración anteriores, es el semibemol el que se emplea para señalar las alteraciones de cuarto de tono, siendo esto suficiente en la casi totalidad de casos. Por la forma particular en la que se crean y se trasponen los modos melódicos como conjuntos de estas 24 notas, no todas ellas se utilizan en la práctica, lo que permite la transcripción de la música árabe con el empleo de una sola alteración adicional.

    No obstante, es importante recalcar que este enfoque no es más que una aproximación para poder escribir la música árabe, y el uso de la microtonalidad en esta es en realidad más complejo. Las notas situadas entre dos notas de la escala cromática no se encuentran en la mitad exacta del intervalo, y su posición ni siquiera es constante. Dependiendo de la época, estilo o región geográfica, existen variaciones, de manera que un semibemol puede situarse más cerca de la nota natural o de la nota bemol, según sea las circunstancias. Incluso si estas son iguales o se trata de un mismo intérprete, la nota puede desplazarse según se toque en un modo melódico u otro.

    A efectos de este libro, obviaremos esta circunstancia y asumiremos la estabilidad de las notas que utilizan intervalos de cuarto de tono, así como de las restantes. Los valores de las tablaturas serán de $\frac12$, y las notas se situarán en los esquemas del diapasón en el punto medio entre las posiciones de las notas adyacentes, aun cuando el lugar en el que se ha de pulsar la nota no ha de ser necesariamente ese.

    Aunque se trata de una terminología en desuso hoy en día, existe una nomeclatura tradicional para las notas musicales. Se muestran a continuación algunos de sus nombres:

\vspace*{0.5cm}
\begin{center}
\begin{lilypond}
\include "arabic.ly"

\paper {
  #(define fonts
    (set-global-fonts
     #:roman "Linux Libertine O,serif"
     #:sans "Linux Biolinum O,sans-serif"
     #:typewriter "Linux Libertine Mono O,monospace"
   ))
}

music =  \relative do {
  \override Staff.TimeSignature #'stencil = ##f
  \override Score.RehearsalMark #'self-alignment-X = #LEFT
  \time 32/4
sol1 la sib do re mib misb mi fa fad sol lab sib sisb do re
}

text = \lyricmode {
\markup {\center-align \fontsize #-2 "Yakah  "} 1
\markup {\center-align \fontsize #-2 "Ushayran  "} 1
\markup {\center-align \fontsize #-2 "qarar Ayam  "}1 
\markup {\center-align \fontsize #-2 "Rast  "} 1
\markup {\center-align \fontsize #-2 "Dukah  "} 1
\markup {\center-align \fontsize #-2 "Kurd  "} 1
\markup {\center-align \fontsize #-2 "Sikah  "} 1
\markup {\center-align \fontsize #-2 "Buselik  "}1
\markup {\center-align \fontsize #-2 "Yiharkah  "}1
\markup {\center-align \fontsize #-2 "Hiyaz"}1
\markup {\center-align \fontsize #-2 "Nawa"}1
\markup {\center-align \fontsize #-2 "Hisar"}1
\markup {\center-align \fontsize #-2 "Ayam"}1
\markup {\center-align \fontsize #-2 "Awk"}1
\markup {\center-align \fontsize #-2 "Kirdan"}1
\markup {\center-align \fontsize #-2 "Muhayyer"}1
}

\score {
  <<
  \new Lyrics \with {
    \override LyricText #'self-alignment-X = #LEFT
    \override LyricText #'font-size = #.2
  }
  { \text }
  \new StaffGroup <<
    \new Staff <<
      \context Voice { \clef "G_8" \music }
    >>
  >>
  >>
  \layout {
    indent = #0
     \context {
         \Score
         \remove "Bar_number_engraver"
     }
  }
}
\end{lilypond}
\end{center}

\vspace*{0.5cm}

Como puede verse, lo anterior no cubre todo el rango del ud, y hay notas que pueden tocarse en este (y también en otros instrumentos) para las que no existen un nombre tradicional. Esto es así porque, una vez más, la voz es el instrumento fundamental de la música árabe, y tan solo se han definido nombres para aquellas notas que representan en el rango vocal típico de un cantante. A pesar de que las notas fuera de esa extensión son posibles con el ud (o incluso puede serlo con la voz, según el intérprete), a efectos prácticos se considera que la melodía ha de estar en esa extensión de notas definidas y con nombres particulares.

Algunos de estos nombres volverán a aparecer cuando veamos los conceptos de \yins y \maqamb, ya que tanto unos como otros toman sus nombres de los de las notas a partir de las que comienzan.

Es relevante también el hecho de que los nombres no se repiten al llegar a la octava, y la misma nota en octavas distintas se considera, según este enfoque tradicional, como una nota diferente.

En la actualidad, como se ha dicho, no se emplean estos nombres, y en su lugar se utiliza la misma nomenclatura que en idioma español (esto es, \emph{do, re, mi\ldots}). La notación por medio de letras (A, B, C\ldots), habitual en países de habla inglesa, así como en otros contextos musicales, no se utiliza en la música árabe.

El sistema tonal árabe, en conjunto con la manera de utilizar sus elementos, conforman lo que se conoce como el \maqam árabe. Por su importancia, lo trataremos de manera independiente en el capítulo \ref{maqam}.

\label{instrumentos}
\section{Instrumentos}

    La música árabe emplea numerosos instrumentos, y en la actualidad a estos se suman la casi totalidad de los occidentales. Sin embargo, en su forma tradicional esta colección se puede reducir a seis de ellos: cuatro melódicos y dos de percusión. Estos instrumentos forman lo que se conoce como el \emph{tajt} o conjunto clásico.

    \subsection{Instrumentos melódicos}

    Además del propio ud, el \emph{tajt} incluye el \emph{ney} o flauta árabe, el \emph{qanun} y el violín.

    \subsubsection{\emph{Ney}}

    El \emph{ney} es una flauta hecha de caña, con seis agujeros en la parte superior y uno en la inferior, que se tapa con el dedo pulgar. A pesar de su simplicidad aparente, se trata de un instrumento difícil de tocar, ya que requiere una compleja técnica de embocadura.

    Flautas de diferentes longitudes producen sonidos de distinta altura, y lo habitual es que el músico de \emph{ney} disponga de una colección de ellas de diferentes tamaños, que utilizará según la tonalidad  y el modo melódico de la pieza a tocar.

    \audioclip{taqsimney.mp3}{\emph{Ney}}

    \subsubsection{\emph{Qanun}}

    El \emph{qanun} es un instrumento de cuerda pulsada, de la familia del salterio. Tiene forma trapezoidal, y sus cuerdas se agrupan en grupos de tres de ellas con el mismo sonido, salvo las más graves, que pueden estar en grupos de dos o ser cuerdas simples. Se toca plano sobre las piernas o apoyado sobre una mesa, pulsando las cuerdas mediante púas que se enganchan a los dedos.

    En su parte izquierda, tiene unas llaves fabricadas en latón que sirven para variar la altura de la nota dada por cada cuerda, y mediante las cuales se pueden obtener los intervalos de cuarto de tono.

    \audioclip{taqsimqanun.mp3}{\emph{Qanun}}

    \subsubsection{Violín}

    El violín árabe no se diferencia del violín occidental. Originalmente se empleaban instrumentos tradicionales de cuerda frotada, tales como el violín egipcio de dos cuerdas, pero fueron reemplazados por el violín en la segunda mitad del siglo \textsc{XIX}.

    La afinación más extendida en la música árabe es sol-re-sol-re (de la cuerda más grave a la más aguda), distinta de la afinación sol-re-la-mi de uso generalizado en el violín europeo.

    Aunque suele tocarse en la posición habitual, sujeto entre el hombro y la barbilla, la postura clásica árabe es distinta: el violín se sitúa en vertical, con su base apoyada sobre el muslo o la rodilla del interprete. En esta posición, para cambiar de cuerdas no se varía el ángulo del arco, sino que se rota el violín sobre su apoyo en la pierna.

    \audioclip{taqsimviolin.mp3}{Violín}

    \subsection{Instrumentos de percusión}

        El instrumento de percusión principal en la música árabe tradicional es el \emph{riqq}. A él se ha sumado, a partir de la segunda mitad del siglo \textsc{XX}, el tambor de copa o \emph{derbake}.

        \subsubsection{\emph{Riqq}}

        El \emph{riqq} es una pandereta de unos 20 centímetros de diámetro, con marco de madera y 10 pares de sonajas situadas de manera regular a lo largo de este.

        Es el instrumento de percusión principal dentro de la música árabe, tanto la clásica como la folclórica, y puede producir una gran variedad de sonidos.

        En la música cantada, es habitual que sea el propio cantante quien se acompaña a sí mismo con el \emph{riqq}.

        Con la incorporación del \emph{derbake}, la técnica del \emph{riqq} experimentó una evolución para adaptar su sonoridad y presencia. En particular, el sonido de las sonajas tiene en la actualidad más importancia que el del parche, algo que previamente no sucedía.

        \audioclip{riqq.mp3}{\emph{Riqq}}

        \subsubsection{\emph{Derbake}}

        El \emph{derbake} (también conocido como \emph{tabla} o \emph{darbuka}, entre otros nombres) es un instrumento de la familia de los tambores de copa. Originalmente se fabricaba en barro y con un parche de piel, pero en la actualidad es más habitual el uso de parches de plástico y cuerpo de fibra o metal.

        Se toca o bien sujetándolo bajo el brazo o bien sobre las piernas. Lo habitual es tocarlo con las manos, pero en algunos lugares y tradiciones se toca con los dedos de una mano y con un pequeño palo en la otra a modo de baqueta.

        \audioclip{derbake.mp3}{\emph{Derbake}}


\section{Ritmos}

La música árabe es muy rica en sus estructuras rítmicas. Existe una amplia colección de patrones rítmicos, conocidos como \emph{iqaat} (plural de \emph{iqa}), que se utilizan para crear formas musicales más complejas, con estructuras basadas en uno o varios de estos patrones.

Los instrumentos básicos de la percusión árabe, según ya hemos visto, son el \emph{riqq} y el \emph{derbake}. Es en función de este último que se acostumbran a definir los patrones rítmicos fundamentales, empleando para ello los dos sonidos básicos que produce el instrumento: un sonido grave obtenido al golpear el centro del parche del tambor y un sonido agudo obtenido al golpear en el borde. Para el primero de ellos se usa la onomatopeya \emph{dum} (y se denota con una letra \emph{D}), mientras que para el segundo se emplea \emph{tak} (y se denota con una letra \emph{T}). A estos se le puede añadir un tercer sonido, \emph{ka} (\emph{K}), que representa el golpe sobre el borde con la mano opuesta a la empleada para el sonido \emph{tak}, con menos intensidad.

Es habitual que en una pieza se den modulaciones de un ritmo a otro de similar estructura, sustituyendo por espacio de unos compases el uno por el otro, para así crear un efecto de contraste.

Se recogen a continuación algunos de los \emph{iqaat} más importantes. Las que se presentan aquí son tan solo las formas fundamentales de los ritmos, que se enriquecen después en la interpretación de la misma manera que una melodía puede ornamentarse. Como ejemplo de esto, véase la estructura básica del ritmo conocido como \emph{maqsum}, recogida en la primera línea, y versiones más complejas de este mismo ritmo en las líneas restantes.

\begin{center}
\begin{lilypond}
    <<
      \new DrumStaff \with{
        drumStyleTable = #percussion-style
        \override StaffSymbol.line-count = #1
      }

      \drummode{
        bd8-D bd4-T bd8-T bd4-D bd4-T
      }
    >>
    <<
      \new DrumStaff \with{
        drumStyleTable = #percussion-style
        \override StaffSymbol.line-count = #1
      }

      \drummode{
        bd8-D bd8-T bd8-K bd8-T bd8-D bd8-K bd8-T bd16-T [bd16-K]
      }
    >>
    <<
      \new DrumStaff \with{
        drumStyleTable = #percussion-style
        \override StaffSymbol.line-count = #1
      }

      \drummode{
        bd8-D bd8-D bd16-K bd16-K bd8-T bd8-D bd16-K bd16-K bd8-T bd16-K bd16-K
      }
    >>
\end{lilypond}
\end{center}

\audioclip{ritmoornamentos.mp3}{Ritmo básico y ritmo con ornamentos}

\subsubsection{Ayub}

Ritmo utilizado en la música sufí y en la música popular árabe. Se toca generalmente a un \emph{tempo} elevado.

\begin{center}
\begin{lilypond}
    <<
      \new DrumStaff \with{
        drumStyleTable = #percussion-style
        \override StaffSymbol.line-count = #1
      }

      \drummode{
      \time 2/4
        bd4-D bd8-D bd4-T
      }
    >>
\end{lilypond}
\end{center}

\audioclip{ritmoayub.mp3}{Ritmo \emph{ayub}}

\subsubsection{Malfuf}

Uno de los ritmos más extendidos, especialmente en la música folclórica. Es habitual modular desde este ritmo al \emph{ayub}, así como al \emph{maqsum}.

\begin{center}
\begin{lilypond}
    <<
      \new DrumStaff \with{
        drumStyleTable = #percussion-style
        \override StaffSymbol.line-count = #1
      }

      \drummode{
        \time 2/4
        bd8-D r16 bd8-T r16 bd8-T
      }
    >>
\end{lilypond}
\end{center}

\audioclip{ritmomalfuf.mp3}{Ritmo \emph{malfuf}}

\subsubsection{Fox}

Ritmo de origen occidental (derivado del \emph{foxtrot}, de donde toma su nombre), se utiliza en la forma musical \emph{longa}. Se toca rápido y con poca ornamentación, dando sensación de vivacidad.

\begin{center}
\begin{lilypond}
    <<
      \new DrumStaff \with{
        drumStyleTable = #percussion-style
        \override StaffSymbol.line-count = #1
      }

      \drummode{
        \time 2/4
        bd8-D bd8-T bd8-D bd8-T
      }
    >>
\end{lilypond}
\end{center}

\audioclip{ritmofox.mp3}{Ritmo \emph{fox}}

\subsubsection{Samai Darij}

Se emplea en la última sección de las que componen la forma musical \emph{samai}, así como en los \emph{darij}.

\begin{center}
\begin{lilypond}
    <<
      \new DrumStaff \with{
        drumStyleTable = #percussion-style
        \override StaffSymbol.line-count = #1
      }

      \drummode{
        \time 3/4
        bd4-D bd4-T bd4-T
      }
    >>
\end{lilypond}
\end{center}

\audioclip{ritmosamaidarij.mp3}{Ritmo \emph{samai darij}}

\subsubsection{Masmudi Saghir}

También conocido como \emph{baladi}, palabra que quiere decir <<del campo>> o <<del pueblo>>, ya que este ritmo proviene de las danzas folclóricas y músicas tradicionales. Existe un estilo de música denominado \emph{baladi}, que aparece a principios del siglo \textsc{XX}, y que puede llevar a confusión. No son conceptos equivalentes, ya que en este estilo aparecen otros ritmos además del \emph{masmudi saguir}, y este a su vez se emplea en formas musicales distintas.

Su estructura es la siguiente:

\begin{center}
\begin{lilypond}
    <<
      \new DrumStaff \with{
        drumStyleTable = #percussion-style
        \override StaffSymbol.line-count = #1
      }

      \drummode{
        bd8-D bd4-D bd8-T bd4-D bd4-T
      }
    >>
\end{lilypond}
\end{center}

\audioclip{ritmomasmudisaguir.mp3}{Ritmo \emph{masmudi saguir}}

\subsubsection{Maqsum}

El más extendido y reconocible de los ritmos árabes. Modula habitualmente a \emph{malfuf}, \emph{baladi} y \emph{saidi}.

\begin{center}
\begin{lilypond}
    <<
      \new DrumStaff \with{
        drumStyleTable = #percussion-style
        \override StaffSymbol.line-count = #1
      }

      \drummode{
        bd8-D bd4-T bd8-T bd4-D bd4-T
      }
    >>
\end{lilypond}
\end{center}

\audioclip{ritmomaqsum.mp3}{Ritmo \emph{maqsum}}

\subsubsection{Saidi}

\begin{center}
\begin{lilypond}
    <<
      \new DrumStaff \with{
        drumStyleTable = #percussion-style
        \override StaffSymbol.line-count = #1
      }

      \drummode{
        bd8-D bd8-T r4 bd8-D bd8-D r8 bd8-T r8
      }
    >>
\end{lilypond}
\end{center}

\audioclip{ritmosaidi.mp3}{Ritmo \emph{saidi}}

\subsubsection{Wahda}

El ritmo más habitual en la música vocal.

\begin{center}
\begin{lilypond}
    <<
      \new DrumStaff \with{
        drumStyleTable = #percussion-style
        \override StaffSymbol.line-count = #1
      }

      \drummode{
        bd8-D r8 r8 bd8-T r8 r8 bd8-T r8
      }
    >>
\end{lilypond}
\end{center}

\audioclip{ritmowahda.mp3}{Ritmo \emph{wahda}}


\subsubsection{Masmudi Kabir}

Ritmo de origen egipcio, se utiliza también en la música de danza. Tiene dos variantes principales: una con dos golpes \emph{Dum} y otra con tres:

\begin{center}
\begin{lilypond}
    <<
      \new DrumStaff \with{
        drumStyleTable = #percussion-style
        \override StaffSymbol.line-count = #1
      }

      \drummode{
      \time 8/4
        bd4-D bd4-D r4 bd4-T bd4-D r4 bd4-T bd4-T
      }
    >>
    <<
      \new DrumStaff \with{
        drumStyleTable = #percussion-style
        \override StaffSymbol.line-count = #1
      }

      \drummode{
      \time 8/4
        bd4-D bd4-D bd4-D bd4-T bd4-D r4 bd4-T bd4-T
      }
    >>
\end{lilypond}
\end{center}

\audioclip{ritmomasmudikabir.mp3}{Ritmo \emph{masmudi kabir}}

\subsubsection{Chiftetelli}

Ritmo de origen turco y uno de los más utilizados en la actualidad en la danza árabe.

\begin{center}
\begin{lilypond}
    <<
      \new DrumStaff \with{
        drumStyleTable = #percussion-style
        \override StaffSymbol.line-count = #1
      }

      \drummode{
      \time 8/4
        bd4-D bd8-D bd4-T bd8-D bd4-D bd4-D bd4-D bd4-T
      }
    >>
\end{lilypond}
\end{center}

\audioclip{ritmochiftetelli.mp3}{Ritmo \emph{chiftetelli}}

\subsubsection{Sombati}

Considerado por algunos como la versión egipcia del ritmo \emph{chiftetelli}.

\begin{center}
\begin{lilypond}
    <<
      \new DrumStaff \with{
        drumStyleTable = #percussion-style
        \override StaffSymbol.line-count = #1
      }

      \drummode{
      \time 8/4
        bd4.-D bd4.-T bd4-T bd4-D r4 bd4-T r4 
      }
    >>
\end{lilypond}
\end{center}

\audioclip{ritmochiftetelli.mp3}{Ritmo \emph{sombati}}


\subsubsection{Samai Thaqil}

Ritmo de diez tiempos divididos en una estructura $3+2+2+3$. Se utiliza en la forma musical \emph{samai} y en las denominadas \emph{moaxajas}.

\begin{center}
\begin{lilypond}
    <<
      \new DrumStaff \with{
        drumStyleTable = #percussion-style
        \override StaffSymbol.line-count = #1
      }

      \drummode{
        \time 10/8
        bd8-D r8 r8 bd8-T r8 bd8-D bd8-D bd8-T r8 r8
      }
    >>
\end{lilypond}
\end{center}

\audioclip{ritmosamaithaqil.mp3}{Ritmo \emph{samai thaqil}}

\subsubsection{Zarafat}

Ritmo popular empleado en las \emph{moaxajas}.

\begin{center}
\begin{lilypond}
    <<
      \new DrumStaff \with{
        drumStyleTable = #percussion-style
        \override StaffSymbol.line-count = #1
      }

      \drummode{
        \time 10/8
        bd8-D r8 r8 bd8-T r8 r8 bd8-D r8 bd8-D bd8-D bd8-T r8 r8
      }
    >>
\end{lilypond}
\end{center}

\audioclip{ritmozarafat.mp3}{Ritmo \emph{zarafat}}

\section{Formas musicales}

    Para el interprete de ud, es necesario conocer todas las formas musicales de la música árabe, tanto vocales como instrumentales, ya que el ud participa de todas ellas. En las instrumentales, por ser el elemento solista por excelencia. En las vocales, por ser también instrumento de acompañamiento, y porque estas acostumbran a incluir pasajes solistas donde el ud es protagonista.

\subsection{Formas vocales}

    Las formas vocales de la música árabe son en general compuestas. Es decir, con partes principales vocales y partes instrumentales. Estas últimas aparecen tanto a modo de introducción como en el desarrollo de la pieza, con función de interludio. En ambos casos, la improvisación es elemento habitual, y esos fragmentos solistas suelen incluir una parte importante de ella.

    Las principales formas vocales son las siguientes:

    \subsubsection{Moaxaja}

    La \emph{moaxaja} es una forma poético-musical surgida en Al-Ándalus a partir de la forma literaria del mismo nombre. Se trata de un canto estructurado en estrofas y estribillos, con una métrica (ritmo \emph{samai thaqil}) y rima definidas. 

    A pesar de sus orígenes, la mayor parte de las moaxajas que se interpretan en la actualidad son recientes, del siglo \textsc{xix} y principios del \textsc{xx}, ya que es un género que ha mantenido una cierta popularidad entre los compositores, y en especial en la región de Aleppo.

    \audioclip{lammabada.mp3}{Moaxaja: Lamma Bada Yatathanna}

    \subsubsection{}
    ***

\subsection{Formas instrumentales}

    Siendo la música árabe eminentemente vocal, las formas musicales instrumentales no han adquirido plena entidad hasta épocas recientes. Es a partir de comienzos del siglo \textsc{XX} que la música instrumental alcanza suficiente categoría como para interpretarse por sí misma. Antes de ello, las piezas instrumentales aparecían tan solo como parte de piezas vocales, o bien se interpretaban intercaladas entre aquellas dentro de conciertos que se articulaban alrededor de la labor de un cantante.

    Distinguimos aquí las formas de origen otomano (\emph{bashraf}, \emph{samai} y \emph{longa}), que se han adaptado y han encontrado su lugar dentro del contexto árabe, y las que se originan directamente dentro de este.

    \subsubsection{Bashraf}

        Proviene del \emph{Peshrev} turco. Tiene una estructura dividida en cuatro o cinco partes denominadas \emph{janas}, entre las cuales se intercala un \emph{taslim} o estribillo.

        La primera \emph{jana} y el \emph{taslim} emplean el mismo \maqamb, pero en las restantes \emph{janas} se dan modulaciones.

        La pieza se cierra siempre con el \emph{taslim}, regresándose así al \maqam original.

        Suele emplear ritmos complejos, que no varían a lo largo de la pieza.

        \audioclip{bashraffarahfaza.mp3}{Bashraf: Bashraf Farahfaza}

    \subsubsection{Samai}

        Al igual que el \emph{bashraf}, el \emph{samai} se compone de cuatro de \emph{janas} entre las que se toca un estribillo. La diferencia principal reside en el ritmo, ya que en el caso del \emph{samai} se emplea el \emph{samai zaqil} para las tres primeras \emph{janas} y un ritmo en $\frac{3}{4}$ para la última (habitualmente \emph{samai darij}).

        \audioclip{samaialaryan.mp3}{Samai: Samai Bayati Al-Aryan}

    \subsubsection{Longa}

        La \emph{longa} comparte con las anteriores formas la estructura de \emph{janas} y \emph{taslim}, aunque esta es en este caso menos estricta y permite variaciones. Es una música de danza que se toca generalmente a velocidad rápida, y que utiliza el ritmo \emph{fox} en $\frac{2}{4}$.

        \audioclip{longafarahfaza.mp3}{Longa: Longa Farahfaza}

    \subsubsection{Tahmila}

        La \emph{tahmila} es una forma musical de origen egipcio. Se toca habitualmente por un \emph{tajt}, y tras una introducción o encabezamiento, todos los miembros de la formación van sucesivamente teniendo turnos para interpretar un solo basado en la melodía principal. Tras cada solo, el grupo al completo toca un estribillo breve, de forma que la \emph{tahmila} tiene un patrón de llamada-respuesta muy evidente.

        Se trata una música enérgica empleada para la danza, y como norma general suele emplear el ritmo \emph{maqsum}.

        \audioclip{tahmilasuznak.mp3}{Tahmila: Tahmila Suznak}

    \subsubsection{Dulab}

        Pieza instrumental breve que se utiliza para presentar el modo melódico o \maqam de una pieza mayor. Contiene todas las características que permiten identificar el \maqam, que luego se desarrolla a lo largo de la composición que sigue.

        \audioclip{dulabrast.mp3}{Dulab: Dulab Rast}

    \subsubsection{Darij}

        Composición instrumental también breve, generalmente sin modulaciones, y que utiliza el ritmo \emph{samai darij}. 

        En la actualidad se interpretan con poca frecuencia, y es raro que se compongan nuevas piezas que entren en esta categoría. Aun así las composiciones existentes, en su mayoría anteriores al siglo \textsc{xx}, son de interés como piezas de estudio de un determinado \maqam, como sucede en este caso.

        \audioclip{darijbayati.mp3}{Darij: Darij Bayati Al Shawwa}

    \subsubsection{Maqtuah}

        Denominación genérica para las piezas instrumentales, sin un determinado ritmo o estructura. Suele estar formada por diversas secciones, entre las cuales puede haber partes improvisadas.

        Muchas de las composiciones que entran dentro de esta categoría se utilizan para el baile, y se denominan \emph{raqsat} (danza).

        \audioclip{aziza.mp3}{Maqtuah: Aziza}

    \subsubsection{Taqsim}

        El \taqsim es una pieza improvisada por un único músico, y constituye una de las máximas expresiones de la música árabe. El \emph{taqsim} es una forma ritmicamente libre (aunque existen variantes con acompañamiento rítmico), y se ciñe tan solo al contexto definido por un \maqam, desarrollándolo a lo largo de su duración.

        Un \taqsim puede formar parte de una pieza mayor (por ejemplo, una \emph{maqtuah}), en cuyo caso actúa a la manera de interludio solista o introducción, o bien ser una pieza por sí mismo.

        De los instrumentos de la música árabe, el violín y el ud son los que con más frecuencia interpretan esta clase de improvisaciones, y en el ud en particular son parte imprescindible del repertorio.

        Debido a su importancia, le dedicaremos una sección independiente (ver \ref{taqsim2}) una vez que hayamos visto en detalle el concepto de \maqam, necesario para comprender la estructura de un \taqsim y la manera de tocarlo.

        \audioclip{taqsimbayati.mp3}{Taqsim: Taqsim Bayati}




